\section{Limitations and Open Problems}
\label{sec:limitations}

The theorem line isolates a baseline independence model. The remaining gaps are
structured as follows.

\begin{enumerate}[leftmargin=*,label=(\arabic*)]
\item \textbf{Finite-sample calibration.} The oracle null law for $Q^*$ is exact,
while the feasible statistic relies on a delta-method approximation. A
nonasymptotic bound for the gap $Q-Q^*$ would quantify the finite-$n$ cost of
feasibility.

\item \textbf{Unknown $\sigma_0$.} A pilot plug-in estimator gives a
self-normalized path for inference, and a parametric bootstrap yields stable
calibration on the tested grid. The oracle split-$F$ benchmark isolates a clean
decoupled route, but the fully feasible split-$F$ procedure remains
conservative. The open theorem problem is to connect the benchmark split-$F$
law to a fully feasible procedure, either by proving the validity of a
cross-fitted statistic or by proving the validity of the same-sample bootstrap.

\item \textbf{Cross-lag dependence.} Shared present-time sampling induces
covariance across lags. The diagonal model is conservative, but a covariance-
adjusted analogue remains to be proved.

\item \textbf{Growing-lag and broader departures.} The adaptivity result covers
consecutive lags under the exact power-law family. Extending the lower
and upper theory to more general lag schedules and to departures from exact
power-law geometry is open.
\end{enumerate}
